One has to notice that in the case of the Sun $d = R_\odot$, that is $d$
equals the solar radius. \hfill(1 p)

The angular deflection scales as $m/r$ (mass over radius) of the deflector
body, so one has to consider $M_{\mathrm{b}}/R_{\mathrm{b}}$, where b stands
for Jupiter and Moon, respectively.
\[
\frac{\Delta\theta_{\mathrm{b}}}{\Delta\theta_\odot}
  = \frac{M_{\mathrm{b}}/R_{\mathrm{b}}}{M_\odot/R_\odot}
  \tag*{(2.1)\quad(2 p)}
\]
This means that:
\[
\Delta\theta_{\mathrm{b}}
  = \Delta\theta_\odot\,\frac{M_{\mathrm{b}}/R_{\mathrm{b}}}{M_\odot/R_\odot}
  \tag*{(2.2)\quad(1 p)}
\]
\begin{description}
\item[a)] Jupiter

Substituting the masses and radii one gets:
\[
\Delta\theta_{\mathrm{J}} = \boxed{0.017''} = \boxed{17~\mathrm{mas}}
  \tag*{(2 p)}
\]
The deflection of radio waves by Jupiter's gravitational field is
\emph{greater} than 0.1\,mas, so it is \emph{observable} with the VLBI
network. The answer is: YES \hfill(1 p)

\item[b)] Moon

Substituting the masses and radii one gets:
\[
\Delta\theta_{\mathrm{M}} = \boxed{0.000\,026''} = \boxed{0.026~\mathrm{mas}}
  \tag*{(2 p)}
\]
The deflection of radio waves by Moon's gravitational field is \emph{smaller}
than 0.1\,mas, so it is \emph{not observable} with the VLBI network. The
answer is: NO \hfill(1 p)
\end{description}
